\section{Metrics}

% Lead owner: Keerthi
% IFS framing input: Sreeja

This section should stay short and operational. The paper does not need a large
metric taxonomy; it only needs the metrics required to evaluate whether PBCP
prevents waste and whether its reasoning aligns with observed behavior.

\subsection{Cost Prevention Score}

Introduce CPS as the primary prevention metric. This subsection should answer a
single question: how much potential waste did the system prevent before or during
execution?

\subsection{Execution Success Rate}

Introduce ESR as the anti-gaming constraint. The point is to show that high
prevention is only meaningful when workloads still complete successfully and the
system is not inflating CPS by over-blocking useful work.

\subsection{Intent-Fit Score}

Introduce IFS as the diagnostic metric. This subsection should make the
relationship clear: CPS says how much waste was prevented, while IFS helps
explain why a workload diverged from its expected behavior.

At first mention in the final prose, keep the canonical terminology stable:
Intent-Behavior Discrepancy (IBD) is the systems framing, and Intent-Fit Score
(IFS) is the metric used to measure semantic alignment between intent and
behavior.
